\section{Background}
\label{sec:background}

\subsection{Concurrency in Go}
Go provides lightweight threads called \emph{goroutines}, started with the
\texttt{go} keyword and scheduled by the Go runtime onto operating-system threads.
Goroutines can communicate in two ways:

\begin{itemize}
  \item \textbf{Message passing.} A \emph{channel} carries typed values between
        goroutines. A send on an unbuffered channel blocks until another goroutine
        receives from it; a buffered channel blocks only when its buffer is full.
        The \texttt{select} statement waits on several channel operations at once,
        and a \texttt{default} branch makes it non-blocking. Closing a channel twice
        causes a run-time panic.
  \item \textbf{Shared memory.} The \texttt{sync} package offers \texttt{Mutex},
        \texttt{RWMutex}, \texttt{Cond}, \texttt{Once} and \texttt{WaitGroup}.
        A \texttt{WaitGroup} counts outstanding tasks: \texttt{Add} increases the
        counter, \texttt{Done} decreases it, and \texttt{Wait} blocks until it reaches
        zero. Anonymous functions (closures) started as goroutines can also share
        variables of the enclosing function, including loop variables.
\end{itemize}

\subsection{Bug Taxonomy Used by the Paper}
Tu et al.~\cite{tu2019go} classify concurrency bugs along two dimensions.
By \emph{behavior}, a \textbf{blocking} bug leaves one or more goroutines stuck
forever (for example, a send that no goroutine will ever receive), while a
\textbf{non-blocking} bug lets all goroutines continue but produces a wrong result,
a data race or a panic. By \emph{cause}, a bug stems from misuse of
\textbf{shared memory} or of \textbf{message passing}.
\todo{Verify the exact counts against the paper before quoting them.}

\subsection{Go's Built-in Bug Detectors}
\begin{itemize}
  \item \textbf{Runtime deadlock detector.} When every goroutine is blocked, the
        runtime aborts with \texttt{fatal error: all goroutines are asleep -
        deadlock!}. It does not fire if any goroutine can still run, so a leak of a
        few goroutines in a live program goes unreported.
  \item \textbf{Data race detector.} Building with \texttt{-race} instruments memory
        accesses (based on ThreadSanitizer) and prints \texttt{WARNING: DATA RACE}
        for unsynchronized conflicting accesses that actually happen during the run.
        It cannot report races on paths that were not executed.
\end{itemize}

\subsection{Bugs Reproduced in This Work}
\Cref{tab:cases} lists the four bugs we reproduce, one per figure of the paper.

\begin{table}[h]
  \centering
  \caption{Bug cases reproduced in this work.}
  \label{tab:cases}
  \begin{tabular}{lllll}
    \toprule
    Paper figure & Original bug & Behavior & Cause & Our test \\
    \midrule
    Figure 1  & Kubernetes \#5316 & Blocking     & Message passing & \texttt{k8s\_5316\_test.go} \\
    Figure 10 & Docker \#24007    & Non-blocking & Message passing & \texttt{docker\_24007\_test.go} \\
    Figure 5  & Docker \#25384    & Blocking     & Shared memory   & \texttt{docker\_25384\_test.go} \\
    Figure 8  & Docker \#27037    & Non-blocking & Shared memory   & \texttt{docker\_27037\_test.go} \\
    \bottomrule
  \end{tabular}
\end{table}

\begin{itemize}
  \item \textbf{Kubernetes \#5316 (Figure 1).} A child goroutine sends its result on
        an unbuffered channel, but the parent has already left through a timeout
        branch of \texttt{select}. The send blocks forever and the child goroutine
        leaks.
  \item \textbf{Docker \#24007 (Figure 10).} Two goroutines both take the
        \texttt{default} branch of a \texttt{select} before either closes the channel,
        so the channel is closed twice and the program panics.
  \item \textbf{Docker \#25384 (Figure 5).} \texttt{WaitGroup.Wait()} is called inside
        the loop that starts the tasks, so the parent waits for a task that has not
        been started yet.
  \item \textbf{Docker \#27037 (Figure 8).} Goroutines started in a loop read the loop
        variable \texttt{i} through a closure while the parent keeps updating it, a
        data race that can make several goroutines see the same value.
\end{itemize}
